\chapter{数字产业化、产业数字化与数字贸易}

\begin{examguide}
\textbf{本章考情}："两化"（数字产业化、产业数字化）是 0258 专业课的标配考点（\important{专硕·核心}），企业边界与盈亏平衡的推导是简答题素材，数字化转型的 NPV 决策是计算题素材；数字贸易规则（DEPA、CPTPP、RCEP）与跨境电商的监管安排是论述题与辨析题的高频素材。

\textbf{核心考点}：（1）ICT 产业的构成与"两化"的测算；（2）ICT 贡献的直接与间接分解；（3）"卡脖子"与产业链安全；（4）科斯企业边界、企业边界收缩与科层制扁平化；（5）数字化转型的成本结构效应与盈亏平衡门槛；（6）数字化投资的 NPV 决策与实物期权；（7）工业互联网、产业数字化的形态与中小企业转型的障碍及政策；（8）去中介化与再中介化；（9）幂律分布、长尾与范围经济；（10）数字贸易的两口径辨析；（11）数字贸易规则体系与跨境电商的监管安排；（12）产业政策的理论依据、数字经济三分类与中外产业政策对照。

\textbf{本章主线}：数字产业自身如何增长（11.1 节）→ 数字技术如何改造传统产业的生产组织（11.2 节）→ 数字渠道如何改变需求的结构（11.3 节）→ 数字化如何重塑国际贸易的形态与规则（11.4 节）→ 政府与市场如何分工（11.5 节）。
\end{examguide}

第 10 章的结尾将问题延伸到实体经济：技术将如何重构产业组织与国际贸易？本章沿五条线索回答这一问题：11.1 节考察数字产业自身如何增长，11.2 节考察数字技术如何改造传统产业的生产组织，11.3 节考察数字渠道如何改变需求的结构，11.4 节考察数字化如何重塑国际贸易，11.5 节考察政府与市场在数字产业中的分工。五条线索共享一个分析框架：数字化改变成本结构，成本结构改变组织边界，组织边界的调整最终以企业形态、产业形态、贸易规则与政策安排的形式呈现出来。

\section{数字产业化}

数字产业化与产业数字化是理解数字经济的两个基本维度，其划分依据已在第 1 章给出：\textbf{数字产业化}指数字技术的产业化部门，其产出本身就是数字产品或数字服务；\textbf{产业数字化}指数字技术对传统产业的渗透改造。本章先考察数字产业化：它的构成是什么、它对经济增长的贡献如何测算、它的结构性短板在哪里。

\subsection{ICT 产业的构成}

数字产业化的主体是\textbf{信息通信技术（ICT）产业}（Information and Communications Technology），对应《数字经济及其核心产业统计分类（2021）》的前四类。按产业链位置，ICT 产业由四个子行业构成：\textbf{电子信息制造业}居于产业链上游，生产芯片、显示面板、通信设备与消费电子产品，是数字经济的硬件底座；\textbf{软件与信息技术服务业}居于中游，提供基础软件（操作系统、数据库）、应用软件与云服务，把硬件能力转化为可用的信息处理能力；\textbf{互联网行业}居于下游，以平台、内容与在线服务直接面向最终用户；\textbf{电信业}提供网络运营与接入服务，构成连接一切的物理通道。四个子行业的分工可以概括为：硬件提供算力与设备，软件定义功能，网络传输信息，互联网组织服务。

ICT 产业在整个数字经济中的角色是\textbf{供给侧}：它为其他一切数字化活动提供技术、设备与服务，数字产业化的水平因此决定了产业数字化的天花板。没有成熟的本土软件业，制造业的数字化转型就依赖进口系统；没有自主的芯片产业链，数字基础设施的扩建就受制于人。这一供给侧定位同时意味着数字产业化的全球分工位置具有战略意义：中国电子信息制造业规模长期居全球首位，智能手机、显示面板、通信设备的产量均居世界第一，但在产业链的高附加值环节仍然落后，这一结构特征在 11.1.3 节展开。

\subsection{ICT 产业的经济贡献测算}

数字产业化的宏观地位需要用数据说话。ICT 部门自身的增加值份额并不大，但各界普遍认为它的贡献远大于份额。这一判断需要一个严格的测算框架：增加值核算回答"自身有多大"，投入产出表回答"它拉动了多少"，增长核算回答"它让生产率提高了多少"。

\begin{derivation}{贡献的双层分解}{}
设总产出为 $Y$，ICT 部门产出为 $Y_{ICT}$，中间投入为 $M_{ICT}$。第一步，\textbf{增加值与增加值率}。ICT 部门的增加值
\[
VA_{ICT} = Y_{ICT} - M_{ICT},
\]
其增加值率 $v = VA_{ICT}/Y_{ICT}$ 衡量单位产出中新增价值所占的份额。
第二步，\textbf{直接贡献}。将 GDP 增长分解为 ICT 部门与其他部门的贡献：
\[
\Delta Y = \Delta Y_{ICT} + \Delta Y_{\text{其他}},
\]
直接贡献率为 $\Delta Y_{ICT}/\Delta Y$，它刻画的是 ICT 部门自身增长的拉动作用。
第三步，\textbf{间接贡献}。ICT 作为通用技术渗透其他部门，其间接贡献借助投入产出表测算。设部门 $j$ 对 ICT 部门的\textbf{直接消耗系数}为 $a_{ICT,j} = x_{ICT,j}/Y_j$（生产单位 $j$ 部门产出直接消耗的 ICT 投入），由直接消耗系数矩阵 $A$ 可得\textbf{完全消耗系数}矩阵 $B = (I - A)^{-1} - I$，元素 $b_{ICT,j}$ 包含了经由产业链间接传导的全部 ICT 投入。ICT 对部门 $j$ 的完全拉动力为 $b_{ICT,j} \cdot Y_j$，对各产业求和即 ICT 的全部经济拉动。
第四步，\textbf{生产率贡献}。以增长核算分解总产出增长：
\[
\Delta \ln Y = s_{ICT}\,\Delta \ln K_{ICT} + s_{\text{其他}}\,\Delta \ln K_{\text{其他}} + s_L\,\Delta \ln L + \Delta \ln A,
\]
其中 $s$ 为各要素的产出弹性。ICT 的贡献由两部分组成：ICT 资本深化（$s_{ICT}\Delta\ln K_{ICT}$）与包含 ICT 溢出在内的全要素生产率增长（$\Delta\ln A$）。
\end{derivation}

这一框架给出的结论是\important{数字产业化的宏观贡献远大于其自身的增加值份额}：直接贡献只有几个百分点，但经由产业关联与生产率溢出，完全贡献可以数倍于自身规模。$b_{ICT,j}$ 与 $a_{ICT,j}$ 的差距度量产业关联链的放大效应，增加值率 $v$ 的高低则反映经济体在全球 ICT 分工中的位置。美国的经验测算提供了最清晰的证据：在 1995--2000 年的生产率加速期，信息技术生产部门占 GDP 的份额不足 5\%，但对劳动生产率加速增长的贡献约为三分之一（Oliner--Sichel，2000），差距全部来自间接贡献。\mn{"两化"辨析速记：数字产业化是"生产数字产品"的部门，产业数字化是"使用数字技术"的部门。}中国的数据同样支持这一判断：据中国信息通信研究院测算，2022 年数字经济规模约 50.2 万亿元，占 GDP 约 41.5\%；其中数字产业化约 9.2 万亿元，占 GDP 仅约 7.6\%，产业数字化约 41 万亿元，占约 33.9\%。份额不足一成的数字产业化部门，支撑了占 GDP 三分之一的产业数字化部门，这一"小核心、大外围"的结构正是投入产出框架中放大效应的宏观呈现。

\subsection{数字产业化的结构性短板：产业链安全}

测算框架刻画的是数字产业化的贡献面，它的另一面是结构性短板。中国数字产业化"规模全球领先"与"核心环节受制于人"并存，这一短板的性质是产业链安全约束，比普通的竞争力差距更严峻。

短板集中在 ICT 产业链的上游与基础环节。\textbf{集成电路}：高端芯片的设计工具（EDA 软件，Synopsys、Cadence、Siemens EDA 三家企业占据全球市场的绝大部分）、制造设备（极紫外光刻机由荷兰 ASML 独家供应）与先进制程代工均依赖进口，国产芯片的自给集中在成熟制程与设计环节；\textbf{基础软件}：桌面操作系统（Windows、macOS）、移动操作系统（Android、iOS）与高端工业软件的市场由美国企业主导；\textbf{关键零部件与材料}：部分高端电子材料与测试设备同样依赖进口。按微笑曲线（施振荣，1992）的位置看，中国 ICT 产业的优势集中在曲线的中段（组装制造），两端的高附加值环节（上游的芯片与材料、下游的品牌与操作系统）仍然薄弱。

这一短板的性质是\textbf{供给安全约束}，其严重程度远超单纯的贸易逆差问题：关键中间品的专用性强、供应商高度集中，一旦供应中断，下游产业将直接面临停产，贸易政策的波动足以造成产业链的中断风险。\mn{"卡脖子"的产业经济学表述（论述题模板）：不可替代中间品（专用性强）+ 供应商集中（断供即停产）→ 产业链安全成为战略约束。政策逻辑三件套：大基金（资本投入）、国产替代（技术追赶）、自主标准（规则保障）。}华为的经历是这一约束的完整样本：2019 年 5 月华为被列入美国实体清单，2020 年 5 月与 8 月芯片供应限制两度升级，台积电等代工厂停止供货，华为高端手机业务一度大幅收缩；2023 年 8 月，搭载国产芯片的 Mate 60 系列发布，标志国产替代在先进制程上取得阶段性突破，但高端芯片的全面自主仍有距离。政策层面的应对同样值得掌握：国家集成电路产业投资基金一期（2014 年，约 1387 亿元）、二期（2019 年，约 2041 亿元）、三期（2024 年 5 月成立，注册资本 3440 亿元）持续投入芯片全产业链，数字产业化的政策含义由此从"做大"扩展为"做安全"。

\section{产业数字化与数字化转型}

数字产业化提供工具，产业数字化使用工具。产业数字化的分析从微观组织开始：数字化首先改变企业的组织边界（11.2.1 节），继而改变成本结构（11.2.2 节），这两者共同决定了企业是否投资数字化以及如何决策（11.2.3 节），最后落到各产业的数字化形态（11.2.4 节）。

\subsection{企业边界理论：数字时代的科斯问题}

数字化转型的最深层影响发生在企业组织层面，其理论基准是科斯（Ronald Coase，1937）在《企业的性质》中提出的企业边界理论。科斯的问题在今天依然尖锐：为什么有的交易在企业内部完成、有的交易在市场完成？数字技术全面改变了两种协调方式的成本，这一问题的数字化版本因此具有直接的现实意义。

\begin{derivation}{企业边界的决定}{}
设企业内部组织一单位额外活动的管理成本为 $MC_{org}$，通过市场购买该活动的交易成本为 $MC_{market}$。企业的最优边界由两种协调方式的边际成本比较决定：
\[
MC_{org} = MC_{market},
\]
当 $MC_{market} > MC_{org}$ 时内部化（纳入企业），当 $MC_{market} < MC_{org}$ 时市场购买（外包）。市场交易成本由三部分构成：搜寻成本（发现交易对手）、缔约与监督成本（订立并执行合同）、协调成本（跨组织协同）。
数字技术沿三条渠道系统性地压低 $MC_{market}$。第一步，\textbf{搜寻成本趋零}：搜索引擎与平台匹配使交易对象的发现成本大幅下降；第二步，\textbf{缔约与监督成本下降}：在线合同、电子支付与评价系统使远距离交易的信任成本降低，陌生人之间的交易成为常态；第三步，\textbf{协调成本下降}：协作软件使跨组织分工的沟通成本下降，地理分散的团队可以像同处一室一样协作。
推论：\important{数字技术使市场交易的相对成本下降，企业边界收缩}。企业的部分职能外包、众包，组织形态从"雇佣"转向"调用"外部资源。
\end{derivation}

这一推论的现实对应广泛存在：平台上的个体经营者（网店店主、独立开发者、零工劳动者）以市场身份承担了传统上企业内部雇员的职能，平台经济中的"个体户化"正是边界收缩的组织表现，也是第 9 章零工经济问题的组织根源。制造业的外包浪潮同样如此：企业保留研发与品牌环节，将制造环节外包给代工厂，组织边界收缩到微笑曲线的两端。

\begin{misconception}
\textbf{误区：数字技术只降低市场交易成本，企业边界必然单向收缩。}数字技术同时也在降低组织成本 $MC_{org}$：管理信息系统使大企业的内部协调更高效，数据驱动的管理使组织可以做得更大而不失控。企业边界的实际变化方向取决于两类成本的相对下降幅度：在交易成本下降更快的环节，边界收缩；在组织成本下降更快的环节，边界甚至可能扩张，数字平台巨头的规模正是后者的体现。考试中以"边界收缩"为标准推论，但完整表述应带上述限定。
\end{misconception}

企业边界讨论的是企业与市场之间的外部组织调整，数字化转型还重塑组织变革的内部维度。传统\textbf{科层制（hierarchy）}组织的信息沿"总部--部门--一线"逐级传递，层级越多，传递的失真与时滞越大，决策权集中在远离信息源的顶层。数字化信息系统使信息在组织内实时共享，单个管理者的有效管理跨度扩大，中间管理层级出现冗余，组织由此走向\textbf{扁平化}：管理层级压缩、决策权下沉，一线团队基于实时数据直接决策。\mn{科层制到扁平化（简答模板）：信息传递成本下降 → 管理跨度扩大 → 中间层级压缩 → 决策权下沉。与误区辨析联动：数字技术同时降低交易成本与组织成本，边界收缩与组织扁平化是同一成本结构变化的两个侧面。}这一转变与边界收缩的经济学同源：信息传递成本是组织成本 $MC_{org}$ 的组成部分，数字化压缩信息传递成本，使更大的组织也可以高效协调，组织的规模上限与协调方式同时改变。

\subsection{数字化转型的成本结构效应}

组织边界的调整伴随成本结构的系统变化。数字化转型需要一次性投入，数字化带来的自动化与数据驱动则压低边际成本，这一"固定成本上升、边际成本下降"的组合改变企业的盈亏平衡条件，是理解企业转型行为异质性的钥匙。

\begin{derivation}{成本结构变化与盈亏平衡}{}
设企业数字化转型的一次性投入（系统建设、数据治理、组织再造）为固定成本 $F$，转型前固定成本为 $F_0$；数字化使边际成本从 $c_0$ 降至 $c_1 < c_0$；价格为 $p$。第一步，转型前后的利润函数：
\[
\pi_0 = (p - c_0) q - F_0, \qquad \pi_1 = (p - c_1) q - (F_0 + F).
\]
第二步，盈亏平衡点：
\[
q_0^* = \frac{F_0}{p - c_0}, \qquad q_1^* = \frac{F_0 + F}{p - c_1}.
\]
第三步，比较静态：
\[
\frac{\partial q_1^*}{\partial F} = \frac{1}{p - c_1} > 0, \qquad \frac{\partial q_1^*}{\partial c_1} = \frac{F_0 + F}{(p - c_1)^2} > 0.
\]
固定投入 $F$ 上升推高盈亏平衡点，边际成本 $c_1$ 下降压低盈亏平衡点，两个效应方向相反。第四步，\important{转型改善利润的条件}：
\[
\pi_1 > \pi_0 \iff (c_0 - c_1) q > F \iff q > q_{\text{门槛}} \equiv \frac{F}{c_0 - c_1}.
\]
\end{derivation}

图~\ref{fig:breakeven} 给出了转型前后两条利润线的对比：转型后的利润线斜率更陡（边际成本下降）但截距更低（固定成本上升），盈亏平衡点从 $q_0^*$ 右移到 $q_1^*$。在 $q_0^*$ 与 $q_1^*$ 之间，不转型的企业已经盈利、转型的企业仍在亏损，这一区间即转型的"门槛区间"。

\begin{figure}[htbp]
\centering
\begin{tikzpicture}[>=Stealth]
\draw[->, very thick] (0,-4.0) -- (0,2.6) node[above=4pt, font=\small] {$\pi$};
\draw[->, very thick] (0,0) -- (5.9,0) node[right=4pt, font=\small] {$q$};
\draw[second, very thick] (0,-1.2) -- (5.58,1.5)
	node[pos=0.62, sloped, above=3pt, font=\small, second] {转型前 $\pi_0$};
\draw[main, very thick] (0,-3.6) -- (5.58,1.8)
	node[pos=0.35, sloped, above=3pt, font=\small, main] {转型后 $\pi_1$};
\fill[white, draw=main, very thick] (2.48,0) circle (2.2pt);
\node[below, font=\small] at (2.48,0) {$q_0^*$};
\fill[white, draw=main, very thick] (3.72,0) circle (2.2pt);
\node[below, font=\small] at (3.72,0) {$q_1^*$};
\draw[<->, very thick] (2.48,-1.6) -- (3.72,-1.6);
\node[below=3pt, font=\small] at (3.1,-1.6) {门槛区间};
\node[left, font=\small] at (0,-1.2) {$F_0$};
\node[left, font=\small] at (0,-3.6) {$F_0+F$};
\end{tikzpicture}
\caption{数字化转型前后的利润线对比}
\label{fig:breakeven}
\end{figure}

门槛效应的经济含义是\important{数字化并非普适的降本手段：只有当产量超过门槛（规模足够大）时，数字化转型才改善利润}。小规模企业的数字化转型可能陷入"投入了固定成本、销量未达门槛"的困境，这一机制解释了数字化转型在企业间的显著异质性：大企业转型积极、中小企业观望的差异，根源在于中小企业的产量位于盈亏平衡点的错误一侧，与其对数字化的态度无关。据中国电子技术标准化研究院《中小企业数字化转型分析报告（2021）》（调查约 1.5 万家中小企业），近八成中小企业处于数字化转型的初步探索阶段，约 12\% 处于应用践行阶段，约 9\% 进入深度应用阶段，门槛效应正是这一分布的重要成因。

门槛效应之外，中小企业的转型障碍还有四类。\textbf{观念}：企业主对数字化转型的收益感知弱，转型意愿不足；\textbf{人才}：既懂业务又懂数字技术的复合型人才稀缺；\textbf{资金}：转型的一次性投入对中小企业的现金流构成压力，融资约束放大门槛效应；\textbf{数据}：企业缺乏数据积累与数据治理能力，转型后的数据价值释放有限。\mn{中小企业转型四挑战速记：观念、人才、资金、数据。政策线"上云用数赋智"三步骤：上云（基础设施）、用数（数据驱动）、赋智（智能升级），概括中小企业数字化从低成本到高价值的递进路径。}政策应对沿"上云用数赋智"主线展开：2020 年 4 月国家发展改革委、中央网信办印发《关于推进"上云用数赋智"行动 培育新经济发展实施方案》，以"上云、用数、赋智"三步骤降低中小企业的转型门槛；2022 年 11 月工业和信息化部印发《中小企业数字化转型指南》，为转型提供标准指引。

此外，数字化还改变了成本的动态结构：数据驱动的学习效应使 $c_1$ 随累计产量进一步下降（学习曲线），先转型者获得持续扩大的成本优势，转型决策由此具有了先发优势的战略属性。

\subsection{企业数字化投资的决策框架}

盈亏平衡回答的是"转型是否可行"，投资决策回答的是"转型是否值得"。当数字化转型被看作一项跨期投资时，标准决策框架是净现值。

\begin{derivation}{数字化转型的 NPV 决策}{}
第一步，净现值公式。设初始投资为 $I_0$（系统、设备、培训），第 $t$ 期由数字化带来的增量现金流为 $\Delta CF_t$（降本 $\Delta C_t$ 与增收 $\Delta R_t$ 之和），折现率为 $r$，项目寿命为 $T$：
\[
NPV = -I_0 + \sum_{t=1}^{T} \frac{\Delta CF_t}{(1+r)^t}, \qquad \Delta CF_t = \Delta C_t + \Delta R_t.
\]
决策规则：\important{$NPV > 0$ 时投资}。第二步，实物期权的引入。数字化转型投资具有三个特征：投资不可逆（系统与组织投入一旦付出难以收回）、收益不确定（$\Delta CF_t$ 的实现滞后且依赖市场状态）、投资时点可推迟。设立即投资的价值为 $NPV_{\text{立即投资}}$，推迟到信息揭晓后再决定的价值为 $NPV_{\text{等待}}$，则项目的真实价值为二者中较大者：
\[
\text{项目价值} = \max\Bigl\{ NPV_{\text{立即投资}},\ E[NPV_{\text{等待}}] \Bigr\},
\]
两项之差即\important{实物期权价值}：当不确定性足够大时，即使 $NPV_{\text{立即投资}} > 0$，若 $E[NPV_{\text{等待}}]$ 更高，最优决策仍是等待；分阶段投资（先试点、再推广）把一次性投入拆分为可延期的投入序列，试点投入相当于为"等待信息"的权利支付的期权费。
\end{derivation}

数字化转型的 NPV 分析有三个特殊之处，它们分别对应着实践中最常见的三类决策错误。其一，\textbf{收益的不确定性}：$\Delta CF_t$ 高度不确定，转型成功率因企业而异，忽视不确定性会把高风险的转型项目误判为优质项目。正确的处理是引入\textbf{实物期权（real option）}分析：投资不可逆、收益不确定且时点可推迟，三者的组合使"等待"本身具有价值，传统 NPV 法则是"现在投或不投"的二元决策，实物期权分析引入了第三选项，\important{即使立即投资的 NPV 为正，等待也可能是最优决策}，先试点再推广的渐进路线由此获得理论支持。其二，\textbf{收益的时滞性}：互补性组织投资未完成前 $\Delta CF_t$ 接近于零，这是第 8 章 GPT 时滞的企业版本；短视的 NPV 计算只截取转型初期的低收益，会系统性低估转型价值。其三，\textbf{外部性}：本企业数字化对产业链上下游的溢出（如大企业数字化倒逼供应商数字化）不计入 $\Delta CF_t$，私人 NPV 低于社会价值，这正是数字化转型补贴政策的理论依据。政策实践与理论推论一致：工业和信息化部、财政部 2023 年启动中小企业数字化转型城市试点，首批约 30 个城市获得中央财政支持，补贴瞄准的正是外部性造成的投资不足。

\subsection{产业数字化的形态}

从企业决策回到产业层面，产业数字化的现实形态沿行业展开，数字化深度取决于各行业产品的性质。\textbf{制造业}的数字化以工业互联网为主线：平台连接设备、工序与供应链，数据驱动工艺优化与预测性维护，"智能制造"的本质是以数据流替代经验流的决策方式。\textbf{C2M}（Customer-to-Manufacturer，用户直连制造）是制造业数字化的需求端形态：消费者订单数据直达产线，生产由"先产后销"转为"以销定产"，个性化定制的成本大幅下降。\textbf{农业}的数字化以传感器、无人机与溯源系统为主要载体，农事决策从经验转向数据。\mn{工业数字化四阶段：机械化（工业 1.0）→ 电气化（2.0）→ 信息化（3.0）→ 智能化（4.0，德国 2013 年"工业 4.0"）。农业数字化转型四路径：生产智能化、经营网络化、管理数据化、服务在线化。C2M：用户直连制造，以销定产。}\textbf{服务业}的数字化深度最深：服务的无形性使数字渠道可以替代物理交付，教育、医疗、金融的在线化程度因此显著高于制造环节。三产深度的梯度在数据上十分鲜明：2020 年服务业数字经济占行业增加值比重约 40.7\%，工业约 21.0\%，农业仅约 8.9\%（中国信息通信研究院，2021），\mn{三产数字化渗透率（信通院 2021）：服务业约 40.7\%、工业约 21.0\%、农业约 8.9\%。规律一句话：服务业最深、农业最浅，与产品可数字化程度一致。}数字化渗透率的梯度与行业产品的可数字化程度一致，也决定了产业数字化政策的发力次序。

工业互联网的中国实践提供了完整的案例样本。海尔卡奥斯 COSMOPlat 平台（2017 年上线）以大规模定制模式连接用户需求与制造环节，是工业和信息化部 2019 年公布的首批十大"双跨"（跨行业、跨领域）工业互联网平台之一；树根互联由三一重工孵化，把工程机械的设备联网与预测性维护经验平台化输出；三一重工长沙 18 号厂房于 2022 年 10 月入选世界经济论坛"灯塔工厂"，是全球重工行业首家灯塔工厂。灯塔工厂的总体数量反映了一国制造业数字化的梯队规模：截至 2024 年 10 月全球约 170 家灯塔工厂中，中国约 70 家，占比居全球首位（世界经济论坛，2024）。

平台在这一过程中的角色是"去中介化"与"再中介化"的双重运动：平台替代了传统中间商（去中介化，如电商替代多级批发），同时自身成为新的、规模更大的中介（再中介化，平台掌握流量分配权）。\mn{去中介化与再中介化的辩证（论述题素材）：数字平台消灭了地域性、层级性的传统中介，但将中介职能集中于自身，中介未消失而是集中化了；平台的流量分配权构成新的渠道权力，渠道成本并未消失，只是换了收取者。}直播电商与本地生活平台的兴起同样如此：它们消灭了地域性、层级性的传统渠道，但将匹配职能集中于自身。

\section{长尾经济}

供给侧的数字化改变了生产组织，需求侧的数字化改变了需求结构。当渠道容量约束被解除后，长期被货架忽略的"尾部需求"首次成为可服务的市场，这就是长尾经济。

\subsection{幂律分布与长尾}

长尾理论由克里斯·安德森（Chris Anderson）在 2004 年发表于《连线》杂志的文章中提出，2006 年扩展为同名著作。长尾的统计基础是幂律分布：在许多数字经济场景中，少数头部商品贡献绝大部分销量，而种类繁多的尾部商品各自销量微小，长尾理论的核心主张正是尾部需求的加总可以媲美头部。幂律分布在第 1 章的网络价值讨论中已经出现（连接价值的分布近似齐普夫分布），这里它再次出现在需求分布中。

\begin{definition}[长尾（long tail）]
长尾是指需求分布中销量微小但种类繁多的尾部商品集合，其经济意义在于：当渠道容量约束被数字技术解除后，尾部商品的总需求可与头部商品相比。
\end{definition}

\begin{derivation}{幂律分布的性质}{}
设商品的销量为 $x$，其销量排名为 $k$，二者近似满足幂律关系
\[
x(k) = C k^{-\alpha}, \qquad \alpha > 0,
\]
其中 $C = x(1)$ 为头名商品的销量。$\alpha = 1$ 即\important{齐普夫定律（Zipf's law）}：排名第 $k$ 的商品销量约为头名的 $1/k$。
第一步，\textbf{头部集中}。销量随排名以幂次衰减，极少数头部商品贡献了绝大部分销量。
第二步，\textbf{货架约束下的可触达销量}。传统零售的货架容量为 $K$ 个 SKU，只能库存头部商品，可触达的销量为
\[
V_{\text{货架}} = \sum_{k=1}^{K} C k^{-1} \approx C \ln K \qquad (\alpha = 1).
\]
第三步，\textbf{长尾的总量}。第 $H+1$ 名到第 $K$ 名的尾部销量为
\[
V_{\text{尾}} = \sum_{k=H+1}^{K} C k^{-1} \approx C(\ln K - \ln H),
\]
而头部 $H$ 名的销量为 $V_{\text{头}} \approx C \ln H$。第四步，\important{长尾超过头部的条件}：
\[
V_{\text{尾}} > V_{\text{头}} \iff C(\ln K - \ln H) > C \ln H \iff \ln K > 2\ln H \iff K > H^2.
\]
当 SKU 总数超过头部数量平方时，长尾的总销量超过头部（$\alpha = 1$）。若 $\alpha < 1$（尾部更"肥"），尾部和随 $K$ 呈多项式增长：
\[
V_{\text{尾}} \approx C\,\frac{K^{1-\alpha}}{1-\alpha},
\]
长尾超过头部的门槛更低。
\end{derivation}

图~\ref{fig:longtail} 给出了幂律分布的双对数表示：销量与排名在对数坐标下呈一条斜率为 $-1$ 的直线，左上浅色区为头部（排名前 10 的商品，单件销量 10--100），右下浅色区为长尾（排名 10--100 的商品，单件销量 1--10），两个区域各自的销量总量相等，这是 $\alpha = 1$ 时幂律分布的"对数均匀"性质；虚线标出了传统货架容量 $K$ 的位置，虚线右侧的尾部需求在传统零售中不可触达，数字渠道使其首次成为可服务市场。

\begin{figure}[htbp]
\centering
\begin{tikzpicture}[>=Stealth]
\draw[->, very thick] (0,0) -- (6.9,0) node[right=4pt, font=\small] {销量排名 $k$};
\draw[->, very thick] (0,0) -- (0,6.9) node[above=4pt, font=\small] {销量 $x(k)$};
\fill[main!12] (0,3.2) rectangle (3.2,6.4);
\fill[second!12] (3.2,0) rectangle (6.4,3.2);
\draw[third, very thick] (0,6.4) -- (6.4,0)
	node[pos=0.4, sloped, above=3pt, font=\small, third] {$x(k) = Ck^{-1}$（齐普夫）};
\draw[dashed, structurecolor, very thick] (4.73,0) -- (4.73,1.67);
\node[below, font=\small] at (4.73,0) {$K$};
\draw (3.2,0.06) -- (3.2,-0.06);
\node[below, font=\small] at (3.2,0) {$10$};
\draw (6.4,0.06) -- (6.4,-0.06);
\node[below, font=\small] at (6.4,0) {$100$};
\draw (0.06,3.2) -- (-0.06,3.2);
\node[left, font=\small] at (0,3.2) {$10$};
\draw (0.06,6.4) -- (-0.06,6.4);
\node[left, font=\small] at (0,6.4) {$100$};
\node[below left, font=\small] at (0,0) {$1$};
\end{tikzpicture}
\caption{幂律分布下的长尾（双对数坐标）}
\label{fig:longtail}
\end{figure}

长尾理论的经济学含义：数字渠道将市场边界从"货架容量"扩展为"仓储与履约能力"，被传统零售忽略的尾部需求首次成为可服务市场。安德森以亚马逊的图书业务为证据：亚马逊图书销售额的一半以上来自实体书店不经营的书目（安德森，2004），流媒体与数字内容平台进一步把库存约束降为零。但\important{长尾不是免费的}：尾部商品仓储分散、配送距离长，单件履约成本显著高于头部商品，尾部服务的成本结构构成长尾商业化的真实约束，纯线上长尾与线下头部的分工由此形成。

\subsection{范围经济与平台多品类扩张}

长尾的需求侧机会能否被抓住，取决于供给侧能否低成本地经营大量品类，这引出\textbf{范围经济（economies of scope）}的概念。

\begin{derivation}{范围经济的条件与平台的品类扩张}{}
设企业生产两种产品的成本为 $C(q_1, q_2)$，分别生产的成本为 $C(q_1, 0) + C(0, q_2)$。第一步，\important{范围经济条件}：
\[
C(q_1, q_2) < C(q_1, 0) + C(0, q_2),
\]
即联合生产的成本低于分别生产。第二步，范围经济的强度。用范围经济指数度量：
\[
SC = \frac{C(q_1, 0) + C(0, q_2) - C(q_1, q_2)}{C(q_1, q_2)},
\]
$SC > 0$ 表示存在范围经济，$SC$ 越大，联合生产的成本节约越多。第三步，数字平台的成本结构。设平台经营 $n$ 个品类，其固定成本 $F$ 为各品类共享，品类 $i$ 的边际成本为 $c_i$，联合生产成本为
\[
C(q_1, \dots, q_n) = F + \sum_{i=1}^{n} c_i q_i,
\]
若各品类由独立企业分别经营，每家企业都需复制一套固定投入，分别生产的总成本为
\[
\sum_{i=1}^{n} C(q_i, 0, \dots) = nF + \sum_{i=1}^{n} c_i q_i.
\]
第四步，比较。两式相减，联合生产节约的成本为 $(n-1)F$，代入范围经济指数得
\[
SC = \frac{(n-1)F}{F + \sum_{i=1}^{n} c_i q_i} > 0,
\]
\important{品类数 $n$ 越多，$SC$ 越大}：共享固定成本的摊薄效应随品类扩张而增强。
\end{derivation}

数字平台的固定投入（技术架构、用户基础、数据系统）具有高度的跨品类共享性：平台增加一个品类 SKU 的边际成本主要为撮合与履约，其基础设施成本已在主品类摊薄，因此\important{平台的品类扩张存在系统性范围经济，这是"超级平台"多品类渗透（第 4 章平台包抄）的供给侧基础}。范围经济与网络外部性相互叠加：品类越全，用户越集中；用户越集中，新品类入驻的吸引力越大，需求方与供给方的双重规模经济驱动平台从单一市场向"什么都卖"的超级平台演进。亚马逊 1995 年以图书起家，此后扩张为全品类电商；京东以 3C 产品起家，随后扩张至日用百货，二者的路径都是范围经济与需求方规模经济共同作用的结果。

\section{数字贸易}

前三节的分析都在国内市场，数字技术对国际贸易的改造同样深刻：交易环节与交付环节都被数字化。交易环节的数字化产生了跨境电商，交付环节的数字化产生了可数字交付的服务贸易，围绕二者的规则体系成为国际贸易治理的前沿战场。

\subsection{定义与统计口径}

数字贸易的定义看似宽泛，其统计口径的划分却直接决定测算结果，因此先厘清定义，再辨析口径。

\begin{definition}[数字贸易]
数字贸易是指以数字技术为支撑、以数字平台为载体、以\important{数据流动为纽带}的国际贸易活动，其统计口径包括数字订购贸易与数字交付贸易。
\end{definition}

数字贸易横跨货物贸易与服务贸易两个传统统计框架，不对应其中任何单一类别，因此口径的区分是本章\important{辨析的核心}。\textbf{数字订购贸易}：通过数字渠道下单、实物交付的跨境交易，即跨境电商，其贸易对象仍是实物商品，数字化的只是交易环节；\textbf{数字交付贸易}：数字产品与服务的跨境交付（软件、云服务、数字内容、在线教育），贸易对象本身即数字化，交付不经过物理海关。\mn{数字贸易两口径辨析：数字订购 = 数字化交易 + 实物交付（跨境电商，货物贸易框架）；数字交付 = 数字化交易 + 数字交付（软件、云服务，服务贸易框架）。辨析题常考：跨境电商出口的是"货"，统计在货物贸易；云服务出口的是"服务"，统计在服务贸易。}前者统计于货物贸易框架，后者统计于服务贸易框架，数字贸易统计难点的根源即在于它横跨两个框架。

两个口径的规模都值得掌握。跨境电商方面，据海关总署统计，2023 年跨境电商进出口总额约 2.38 万亿元，同比增长约 15.6\%（海关总署，2024），显著快于货物贸易整体增速。数字交付方面，据世界贸易组织统计，2022 年全球可数字化交付服务出口约 3.8 万亿美元，约占全球服务出口总额的 54\%（世界贸易组织，2023），数字交付已构成服务贸易的主体形态。

\subsection{数字贸易的比较优势}

数字贸易不仅改变了贸易的形态，也改变了比较优势的来源。李嘉图（David Ricardo）模型的结论是：两国按相对生产率分工，各自出口相对成本较低的产品。数字技术对这一框架的改造体现在三个层面，其中数据禀赋的引入可以严格刻画。

\begin{derivation}{比较优势的数字化扩展}{}
设本国（$H$）与外国（$F$）生产两种数字服务 $1$、$2$，单位劳动需求分别为 $a_1$、$a_2$ 与 $a_1^*$、$a_2^*$。第一步，李嘉图基准。本国出口服务 $1$ 的条件是
\[
\frac{a_1}{a_2} < \frac{a_1^*}{a_2^*},
\]
即本国在服务 $1$ 上具有比较优势，分工模式由相对生产率的跨国比较决定。
第二步，数据进入生产函数。设服务 $i$ 的单位劳动需求随本国的数据存量 $D$ 下降：
\[
a_i = a_i(D), \qquad a_i'(D) < 0,
\]
设服务 $1$ 相对服务 $2$ 更依赖数据（数据密集度高）。第三步，数据禀赋对比较优势的边际影响。对相对单位劳动需求取对数求导：
\[
\frac{d\ln(a_1/a_2)}{d\ln D} = \frac{D a_1'(D)}{a_1} - \frac{D a_2'(D)}{a_2} < 0,
\]
不等号成立的原因是服务 $1$ 的数据弹性更大。第四步，结论：本国数据存量 $D$ 的上升使 $a_1/a_2$ 下降，比较优势向数据密集型服务移动。\important{数据禀赋的差异成为贸易模式的独立决定因素}，数据大国在数据密集型数字服务上获得传统要素禀赋之外的新优势。
\end{derivation}

上述推导刻画的是三个层面中的第一个：数据禀赋成为比较优势的独立来源。另外两个层面的改造同样重要。其二，\textbf{服务可贸易性革命}：传统服务贸易受"面对面交付"约束，跨境成本极高；数字交付使服务的跨境成本趋零，服务的比较优势首次可以被规模化实现，在线教育、远程医疗、软件即服务（SaaS）的出口都是这一革命的产物。其三，\textbf{长尾的跨境化}：数字渠道使小众产品可以低成本触达全球分散的需求，贸易的规模门槛下降，中小企业的贸易参与度上升，阿里巴巴国际站与速卖通等跨境平台把传统上依赖大型贸易商的出口通道"平台化"，使中小企业直接面向全球消费者，其机制与 11.2.4 节去中介化的逻辑同构。

\subsection{数字贸易规则体系}

数字贸易的快速发展与多边规则的滞后形成鲜明反差，规则体系的建设沿三条轨道推进，构成当前国际贸易治理的前沿战场。\textbf{WTO 框架}：电子商务诸边谈判（联合声明倡议）于 2019 年 1 月达沃斯部长级会议启动，目前约 90 个成员参与，议题集中于跨境数据流动、数据本地化、源代码披露与电子传输关税；电子传输免征关税自 1998 年起延续，2024 年 2 月第 13 届部长级会议同意将豁免延长至 2026 年 3 月。多边进展缓慢的根源在于发达国家与发展中国家的根本分歧：前者主张数据自由流动，后者强调数据主权与产业政策空间。\textbf{区域协定}：DEPA（数字经济伙伴关系协定）由新加坡、新西兰、智利于 2020 年 6 月签署，2021 年 1 月生效，以 16 个模块的模块化结构为特色，成员可选择加入部分模块，中国于 2021 年 11 月正式申请加入；CPTPP 于 2018 年 12 月生效，其电子商务章节（第 14 章）包含数据自由流动、禁止数据本地化与禁止强制披露源代码的高标准条款，中国于 2021 年 9 月正式申请加入；RCEP 于 2022 年 1 月生效，是全球规模最大的自由贸易协定，其电子商务章节（第 12 章）促进无纸化贸易与电子认证，同时设有安全例外条款，对发展中国家留有政策空间。\textbf{国内法外溢}：GDPR 的域外管辖效力使欧盟数据保护标准成为事实上的全球基准，阿努·布拉德福德（Anu Bradford）称之为"布鲁塞尔效应"（第 12 章详述 GDPR 的制度细节）。

三条轨道的推进节奏与约束强度各不相同。\mn{数字贸易规则三集群：WTO 电子商务诸边（多边、缓慢）、DEPA/CPTPP/RCEP（区域、渐进）、国内法域外效力（单边、事实性外溢）。中国立场：申请加入 DEPA 与 CPTPP，支持数据安全有序流动，反对单边数字税与过度数据本地化（第 12 章衔接）。}三个集群的对比见表~\ref{tab:digitaltrade}。

\begin{table}[htbp]
\centering
\setlength{\tabcolsep}{4pt}
\caption{三大区域协定的数字贸易规则对比}
\label{tab:digitaltrade}
\begin{tabular}{p{1.6cm}p{3.2cm}p{3.6cm}p{3.4cm}}
\hline
\textbf{协定} & \textbf{生效与结构} & \textbf{数据流动与本地化} & \textbf{源代码与其他条款} \\
\hline
DEPA & 2021 年 1 月生效；新加坡、新西兰、智利发起；16 个模块，可选择性加入 & 设跨境数据流动与数据创新模块，倡导自由流动 & 含人工智能治理、数字身份等前沿模块 \\
\hline
CPTPP & 2018 年 12 月生效；11 个成员 & 第 14 章：数据自由流动，原则上禁止数据本地化，允许公共政策例外 & 禁止强制披露源代码，高标准 \\
\hline
RCEP & 2022 年 1 月生效；15 个成员，全球规模最大 & 第 12 章：促进数据流动，成员可基于安全例外施加限制 & 无强制源代码条款，对发展中国家留有政策空间 \\
\hline
\end{tabular}
\end{table}

\subsection{跨境数据流动与跨境电商}

规则体系的焦点是跨境数据流动。数据自由流动提升贸易效率：数字交付贸易以数据流动为存在前提，数据本地化要求则构成贸易壁垒。据欧洲国际政治经济中心（ECIPE）测算，数据本地化措施可能使实施国 GDP 下降约 0.1\%--1\%（欧洲国际政治经济中心，2014），并抬高本地服务价格、抑制投资。但数据自由流动同时涉及隐私保护、国家安全与数据主权的多重关切，各国在"自由流动"与"本地化要求"之间形成光谱（第 12 章详述制度细节），数字贸易规则谈判的实质就是在这两端之间寻找折中点。

跨境电商作为数字订购贸易的典型形态，其运作涉及一套专门的监管安排。海关监管代码规定了四种基本模式：\textbf{9610}（跨境电商 B2C 出口"集货报关"，小包裹集货后统一报关出境）、\textbf{9710}（B2B 直接出口，境内企业通过电商平台与境外企业直接交易）、\textbf{9810}（出口海外仓，货物先批量运抵海外仓，境外销售后逐单发货）、\textbf{1210}（保税进口，货物先批量入区，消费者下单后出区清关）。\mn{跨境电商监管代码速记：9610 集货出口（B2C）、9710 直接出口（B2B）、9810 海外仓出口、1210 保税进口。经济逻辑一句话：批量预置 + 逐单清关，用库存风险换时效与成本。}保税仓的运作逻辑是理解跨境电商成本结构的关键：进口商品以批量海运入区（免于逐单报关的高成本），消费者下单后出区时再逐单清关，国内配送时效通常为 1--3 天，显著优于直邮模式的 7--15 天；其经济实质是把"小包化、高频次"贸易的海关成本通过"批量预置 + 逐单清关"摊薄，代价是库存风险由需求预测的精度承担。海外仓（9810）解决了出口方向的同类问题：据商务部统计，截至 2022 年中国企业建设运营的海外仓超过 2000 个（商务部，2022），跨境物流时效从直邮的数周缩短至本地配送的数天。

\section{产业政策：政府与市场的分工}

本章前四节的分析把企业与市场视为资源配置的主体，但数字产业化与产业数字化的推进从来不是纯粹的市场过程：芯片、基础软件、算力网络等领域的投资决策同时取决于企业利润与政府战略。政府应当以什么形式介入、介入到什么程度，是产业分析无法回避的制度问题。本节回答三个递进的问题：产业政策是什么、其理论依据何在、数字经济的产业政策呈现出何种结构与挑战。

\begin{definition}[产业政策（industrial policy）]
产业政策是政府为影响产业结构与产业组织而采取的一系列干预措施。按作用方式分为两类：\textbf{功能性产业政策}（functional industrial policy）通过公共品供给与制度完善改善市场运行条件，如基础设施建设、数据开放、研发税收优惠；\textbf{选择性产业政策}（selective industrial policy）通过补贴、保护、准入安排扶持特定产业或技术路线，如集成电路专项支持。
\end{definition}

产业政策的理论依据有三条。其一，\textbf{市场失灵}：创新与数据流通具有正外部性，私人投资低于社会最优，需要公共投入补足（下文推导）；其二，\textbf{幼稚产业保护}：后发国家的新兴产业在早期成本高于先行者，需要保护期积累规模与经验，再进入国际竞争；其三，\textbf{协调失灵}：产业升级涉及上下游同步投资（设备、材料、工艺的匹配），单个企业无法协调，政府以承诺机制降低协调成本。三条依据在数字经济中都有直接对应：数据要素的外部性、算力与基础软件的追赶、产业链上下游的同步国产化。

\begin{derivation}{研发溢出与最优补贴}{}
\textbf{第一步，私人最优化。}设数字经济核心技术的代表性企业投入研发 $x$，研发的私人收益为 $R(x)$（$R' > 0$，$R'' < 0$），研发成本为凸函数 $C(x)$（$C' > 0$，$C'' > 0$）。企业利润为
\[
\pi(x) = R(x) - C(x),
\]
一阶条件给出私人最优研发 $x^*$：
\[
R'(x^*) = C'(x^*).
\]
\textbf{第二步，社会最优化。}研发具有\textbf{知识溢出}：企业的研发成果以比例 $\beta \in (0,1)$ 外溢，提升其他企业收益 $\beta R(x)$。社会计划者的目标函数为
\[
W(x) = (1+\beta)R(x) - C(x),
\]
一阶条件给出社会最优研发 $x^{**}$：
\[
(1+\beta)R'(x^{**}) = C'(x^{**}).
\]
由于 $R'' < 0$ 且 $C'' > 0$，比较两个一阶条件得 $x^* < x^{**}$：\important{存在知识溢出时，私人研发投资系统性低于社会最优，差距由溢出率 $\beta$ 刻画}。\textbf{第三步，补贴校正。}设政府对每单位研发给予补贴 $s$，企业利润变为 $\pi(x) = R(x) - C(x) + sx$，新的一阶条件为
\[
R'(x^*) + s = C'(x^*).
\]
令 $s = \beta R'(x^{**})$，企业的最优条件退化为社会最优条件：
\[
R'(x^{**}) + \beta R'(x^{**}) = (1+\beta)R'(x^{**}) = C'(x^{**}).
\]
结果：\important{按溢出强度 $\beta R'$ 设定的研发补贴，恰好使私人最优与社会最优重合}，这是研发补贴的规范经济学基础。
\end{derivation}

这一推导把产业政策定位为外部性校正工具的选择问题，市场机制的有效性并不因此被否定。数字经济中的对应十分直接：芯片、基础软件与人工智能基础研究的溢出率高（成果被全产业复用），正是补贴理论的适用场景；而接近消费端的应用创新溢出率低、市场信号清晰，政府介入的理由相应减弱。\mn{产业政策三依据：市场失灵、幼稚产业、协调失灵；数字经济中分别对应数据外部性、算力与软件追赶、上下游同步国产化。}补贴幅度的设定依赖溢出率 $\beta$ 的可度量性：$\beta$ 难以观测时，补贴水平只能近似，甚至被策略性申报扭曲，这正是后文补贴操纵挑战的理论根源。

数字经济的产业政策实践中形成\textbf{三分类}结构。\textbf{功能性政策}构建通用底座：新型基础设施建设（2018 年中央经济工作会议首次提出"新基建"概念）、公共数据开放、全国一体化算力网；\textbf{选择性政策}聚焦战略缺口：国家集成电路产业投资基金（11.1.3 节）、基础软件与工业软件的专项支持；\textbf{竞争性政策}维护市场活力：第 6 章的反垄断与公平竞争审查，防止补贴扶持的冠军企业转身成为垄断者。三分类的分工逻辑是：市场解决应用层，政策解决通用层与瓶颈层，竞争政策守住结构层。

产业政策的经典争论是林毅夫与张维迎之争：产业政策的可行性取决于政府能否在事前识别市场失灵与技术前景，林毅夫以新结构经济学主张按比较优势甄别产业，张维迎则以知识的分散性否定政府甄别能力，主张以保护企业家精神与产权替代补贴。数字经济为这场争论提供了新证据：功能性政策的共识度高（基础设施的外部性确定），选择性政策的风险集中在技术路线赌注（如历史上对平板显示、新能源的不同路线选择），政策设计的重心因此从"事前完美甄别"转向"事后调整机制"。

数字经济产业政策面临的挑战有三。其一，\textbf{信息不对称与补贴操纵}：研发补贴触发策略性申报，企业以低质量研发套取补贴，补贴本身成为企业目标；其二，\textbf{央地协调}：地方政府在数字园区、智算中心上的重复建设与招商补贴竞争，造成算力利用率不足的产能错配；其三，\textbf{与竞争政策的协调}：网络外部性使补贴扶持的市场赢家具有赢者通吃倾向，产业政策制造的集中度需要竞争政策消化。挑战的实质是政府与市场分工的边界问题：\important{数字经济的产业政策有效性取决于三个条件：外部性可度量、技术路线可纠错、竞争政策可制衡}。

数字产业政策的国际对照同样值得掌握。\textbf{美国}以《芯片与科学法案》（2022 年 8 月签署，为半导体制造与研发提供约 527 亿美元补贴）为核心，配合出口管制推动先进制造回流，政策取向是竞争导向与产业回流双轨并行；\textbf{欧盟}发布"2030 数字十年"计划（2021 年），目标包括 2030 年欧洲先进半导体产量占全球 20\%，并出台《欧洲芯片法案》（2023 年，总投入约 430 亿欧元），政策取向侧重规则治理与数字主权；\textbf{日本}提出"社会 5.0"（Society 5.0，2016 年）作为超智能社会的总体愿景，以经济安全保障与半导体制造复兴为政策重点。中外数字产业政策的共同点在于：均将芯片与先进计算列为战略优先、均以补贴与税收优惠为主要工具、均强调数字基础设施的先导作用；路径差异在于：美国以竞争与管制双轨施压，欧盟以规则塑造市场，日本以供应链安全为导向复兴制造，中国则以大基金等举国体制工具布局全产业链自主。这一对照是中外产业政策比较题的答题素材。

\begin{chapterbox}
\textbf{本章考点清单}

\begin{enumerate}
\item ICT 产业的构成与"两化"的测算 \important{【专硕·核心】}
\item ICT 贡献的双层分解：增加值率、直接贡献与投入产出间接贡献 \important{【专硕·核心】}
\item "卡脖子"的产业链安全含义与大基金 \important{【专硕·扩展】}
\item 科斯企业边界 $MC_{org} = MC_{market}$、边界收缩与科层制扁平化（层级压缩、决策权下沉）\important{【专硕·核心】}
\item 数字化转型的盈亏平衡与门槛条件 $q > F/(c_0 - c_1)$ \important{【专硕·核心】}
\item 中小企业转型四挑战（观念、人才、资金、数据）与"上云用数赋智"政策线 \important{【专硕·扩展】}
\item 数字化投资的 NPV 决策与实物期权 \important{【专硕·核心】}
\item 工业互联网、灯塔工厂与产业数字化的行业形态 \important{【专硕·核心】}
\item C2M 用户直连制造、工业数字化四阶段与农业数字化转型四路径 \important{【专硕·扩展】}
\item 去中介化与再中介化的辩证 \important{【专硕·核心】}
\item 幂律分布 $x(k) = Ck^{-\alpha}$、长尾超过头部的条件 $K > H^2$ 与履约约束 \important{【专硕·核心】}
\item 范围经济条件与平台多品类扩张 \important{【专硕·扩展】}
\item 数字贸易两口径（数字订购 vs 数字交付）的辨析 \important{【专硕·核心】}
\item 比较优势的数字化扩展（数据禀赋、服务可贸易性）\important{【专硕·扩展】}
\item 数字贸易规则三集群与 DEPA/CPTPP/RCEP 对比 \important{【专硕·核心】}
\item 跨境电商监管代码与保税仓的运作逻辑 \important{【专硕·扩展】}
\item 产业政策的理论依据（市场失灵、幼稚产业、协调失灵）与研发溢出补贴 $s = \beta R'$ 的推导 \important{【专硕·核心】}
\item 数字经济产业政策的三分类与林张之争 \important{【专硕·扩展】}
\item 中外数字产业政策对照：美国《芯片与科学法案》、欧盟 2030 数字十年与《欧洲芯片法案》、日本社会 5.0 \important{【专硕·扩展】}
\end{enumerate}
\end{chapterbox}

本章的分析贯穿一条主线：数字化改变成本结构，成本结构改变组织与市场的边界，从企业边界到市场边界、再到国家之间的贸易边界，边界调整的方向都由成本变化的方向决定。11.5 节表明，这些边界之外还有一层政府与市场的分工边界，政策工具的合理形态由外部性可否度量、技术路线可否纠错、竞争政策可否制衡三个条件决定。数据该不该跨境流动、隐私如何保护、平台承担何种责任、数字化利润如何征税，则是下一章的制度问题。第 12 章将把数据治理、隐私保护与数字经济税收作为数字经济的制度基础设施展开。
